\usepackage[utf8]{inputenc}

% Fontes e Idioma
\usepackage{mathptmx} % Altera a fonte do texto e da matemática para Times
% O idioma [brazil] já é gerenciado nativamente pelas opções da classe no main.tex

% --- Margens e Espaçamento ---
% O abntex2 já define as margens padrão da ABNT (3cm superior/esquerda, 2cm inferior/direita).
% Não utilize o pacote geometry aqui para não quebrar o layout da classe.

\usepackage{setspace}
\OnehalfSpacing % Espaçamento entre linhas de 1,5

% Elementos Gráficos e Imagens
\usepackage{graphicx} 
\usepackage{caption} % Recursos personalizáveis para legendas
\usepackage{float} 
\usepackage{tikz}
\usepackage{pgfplots} % Plotagem de gráficos
\pgfplotsset{compat=1.18}
\usepackage{adjustbox}

% Tabelas
\usepackage{tabularx}
\usepackage{array}
\usepackage{colortbl}
\usepackage{multicol,multirow}
\setlength{\columnsep}{.75cm}

% Formatação de Texto e Elementos
\usepackage{changepage} % Altera a margem de um bloco
\usepackage{xcolor}
\usepackage{enumitem} % Personalizar listas
\usepackage{pifont} % Símbolos adicionais
\usepackage{textcomp}
\usepackage{textcase} % Deixar títulos em maiúsculo

% --- Configuração de Parágrafos (Normas ABNT) ---
% Removido o pacote parskip que eliminava o recuo do parágrafo.
\setlength{\parindent}{1.25cm} % Recuo padrão da primeira linha do parágrafo
\setlength{\parskip}{0.2cm}    % Espaçamento sutil entre parágrafos

% Matemática Avançada
% (Removido ntheorem pois causava conflito com amsthm)
\usepackage{amssymb,amsmath,amsfonts,amsthm,mathtools,bm,siunitx,empheq,gensymb}
\everymath{\displaystyle}

% --- Elementos Pós-Textuais (Hiperlinks, Citações e Glossário) ---
% Citações padrão ABNT
\usepackage[alf]{abntex2cite}

% Configuração dos links (O abntex2 já carrega o hyperref nativamente)
\hypersetup{
    colorlinks=true,
    linkcolor=black,  % Cor dos links internos (como o sumário)
    citecolor=black,  % Cor dos links para as referências
    urlcolor=black    % Cor dos links de sites
}

% Glossário (Deve ser carregado após o hyperref)
\usepackage[acronym]{glossaries}
\makeglossaries